

\chapter{复围道积分的几何变形与等价性证明}
\label{app:contour}

在复变函数与解析数论中，围道积分的选择往往取决于具体的应用场景与计算便利性。在本书中，我们遇到了两种形式迥异但本质等价的复围道：
\begin{enumerate}
    \item \textbf{锁孔形围道（Hankel 围道，记作 $\mathcal{H}$）}：主要应用于第\ref{ch:gamma_function}章 $\Gamma$ 函数的汉克尔表示以及第\ref{ch:riemann_zeta_continuation}章 $\zeta(s)$ 解析延拓的第一法中。其优势在于能够优雅地绕过复平面的多值分支切割线。
    \item \textbf{矩形围道（记作 $\mathcal{R}$）}：主要应用于第\ref{ch:zero_distribution}章零点个数计数公式 $N(T)$ 的推导以及第\ref{ch:riemann_explicit_formula}章黎曼显式公式中对 Perron 公式的截断误差估计与围道平移。其优势在于几何直观，且垂直边界与实轴平行，易于进行模长估计。
\end{enumerate}

本附录旨在通过严格的拓扑变形（同伦）理论与极限估计，论证这两种围道积分的等价性，从而消除正文各章节之间关于围道选取的逻辑跨越。

---

\section{两种经典围道的几何图像}

首先，我们在复平面中给出两种围道的几何直观表示，如图 \ref{fig:contours} 所示。

\begin{figure}[htbp]
\centering
\resizebox{0.99\textwidth}{!}{%
\begin{tabular}{cc}
\begin{tikzpicture}[>=Stealth, thick]
    % 坐标轴
    \draw[->, thick, gray] (-1.5,0) -- (4.5,0) node[right, text=black] {$\operatorname{Re}(z)$};
    \draw[->, thick, gray] (0,-2.0) -- (0,2.0) node[above, text=black] {$\operatorname{Im}(z)$};
    \def\r{0.7} 
    \def\eps{0.22} 
    \pgfmathsetmacro{\ang}{asin(\eps/\r)} 
    
    % 绘制轮廓
    \draw[blue, thick] (4.0, \eps) -- ({\r*cos(\ang)}, \eps) 
                arc (\ang:360-\ang:\r) 
                -- (4.0, -\eps);
    
    % 方向箭头
    \draw[blue, thick, ->] (2.6, \eps) -- (2.4, \eps); 
    \draw[blue, thick, ->] ({- \r}, 0.02) -- ({- \r}, -0.02); 
    \draw[blue, thick, ->] (2.4, -\eps) -- (2.6, -\eps); 
    
    \node[below left] at (0,0) {$0$};
    \node[blue, above right] at ({1.05*\r*cos(45)}, {1.05*\r*sin(45)}) {$\mathcal{H}$};
    \node[above right] at (1.2, 0.22) {\tiny 上侧：$\arg(-z)=-\pi$};
    \node[below right] at (1.2, -0.22) {\tiny 下侧：$\arg(-z)=\pi$};
    \fill (0,0) circle (1.5pt) node[anchor=south west, yshift=2pt] {\tiny 分支点 $O$};
\end{tikzpicture}
&
\begin{tikzpicture}[
    >=stealth,
    decoration={markings, mark=at position 0.6 with {\arrow[scale=1.3]{>}}}
]
    % 坐标轴
    \draw[->, thick, gray] (-3.5, 0) -- (2.0, 0) node[right, text=black] {$\operatorname{Re}(s)$};
    \draw[->, thick, gray] (0, -2.0) -- (0, 2.0) node[above, text=black] {$\operatorname{Im}(s)$};
    
    % 矩形围道参数
    \def\c{1.0}
    \def\R{-2.8}
    \def\T{1.5}
    
    % 绘制矩形围道
    \draw[thick, red!80!black, postaction={decorate}] (\c, -\T) -- (\c, \T) node[midway, right] {\tiny $\operatorname{Re}(s)=c$};
    \draw[thick, red!80!black, postaction={decorate}] (\c, \T) -- (\R, \T);
    \draw[thick, red!80!black, postaction={decorate}] (\R, \T) -- (\R, -\T) node[midway, left] {\tiny $\operatorname{Re}(s)=-R$};
    \draw[thick, red!80!black, postaction={decorate}] (\R, -\T) -- (\c, -\T);
    
    % 点缀一些极点和零点
    \filldraw[black] (0.8,0) circle (1.2pt) node[below right] {\tiny $1$};
    \filldraw[blue!70!black] (0.3, 0.9) circle (0.8pt) node[right] {\tiny $\rho$};
    \filldraw[blue!70!black] (0.3, -0.9) circle (0.8pt) node[right] {\tiny $\bar{\rho}$};
    \filldraw[green!60!black] (-1.5, 0) circle (1.2pt) node[below] {\tiny $-2$};
    
    \node[red!80!black] at (-1.8, 1.1) {\small 围道 $\mathcal{R}$};
\end{tikzpicture} \\
(a) 锁孔形 (Hankel) 围道 $\mathcal{H}$ & (b) 矩形围道 $\mathcal{R}$ \\
\end{tabular}%
}
\caption{避开正实轴割线的锁孔形围道与包围孤立奇点的矩形围道示意图}
\label{fig:contours}
\end{figure}

---

\section{围道的同伦拓扑变形理论}

两种围道在计算积分时的等价性，在拓扑学上基于\textbf{Cauchy 积分定理（Cauchy's Integral Theorem）}以及\textbf{围道的同伦（Homotopy）变形}。

\begin{theorem}{围道变形原理}{}
设 $D \subset \mathbb{C}$ 是一个单连通区域，$f(z)$ 是定义在 $D \setminus \{z_1, z_2, \dots, z_n\}$ 上的全纯函数，其中 $\{z_k\}$ 为 $f(z)$ 的孤立奇点。若两条正向简单闭曲线 $\mathcal{C}_1$ 与 $\mathcal{C}_2$ 在圆盘 $D$ 内包围了完全相同的奇点集合，则有：
\begin{equation}
    \oint_{\mathcal{C}_1} f(z) dz = \oint_{\mathcal{C}_2} f(z) dz
\end{equation}
\end{theorem}

\begin{proof}
我们可以通过在 $\mathcal{C}_1$ 与 $\mathcal{C}_2$ 之间引入互反向抵消的连接线段，将环形区域（或多连通区域）剖分为若干个不含任何奇点的单连通子区域。根据 Cauchy 积分定理，在每个不含奇点的子区域上，沿边界的积分恒为零。将所有子区域的边界积分相加，内部连接线段上的积分因方向相反相消，最终仅剩下 $\mathcal{C}_1$ 与 $\mathcal{C}_2$ 上的积分（方向相反）。移项后即得两积分相等。
\end{proof}

对于锁孔形围道 $\mathcal{H}$ 与矩形围道 $\mathcal{R}$，只要我们限制在\textbf{切割复平面 $\mathbb{C} \setminus [0,\infty)$} 中，任何被积全纯函数在这两个围道上的行为均符合上述定理的推广。为了将非闭合的锁孔围道与闭合的矩形围道建立联系，我们通常采用\textbf{极限变形（Limit deforming）}的手段，如图 \ref{fig:deformed_contour} 所示。

\begin{figure}[htbp]
\centering
\begin{tikzpicture}[
    >=stealth,
    decoration={markings, mark=at position 0.55 with {\arrow[scale=1.3]{>}}}
]
    % 绘制坐标轴
    \draw[->, thick, gray] (-4.5, 0) -- (3.0, 0) node[right, text=black] {$\operatorname{Re}(z)$};
    \draw[->, thick, gray] (0, -2.5) -- (0, 2.5) node[above, text=black] {$\operatorname{Im}(z)$};
    
    % 分支切割线
    \draw[ultra thick, brown, opacity=0.4, decorate, decoration={zigzag, segment length=1.5mm, amplitude=0.4mm}] (0,0) -- (2.5,0);
    
    % 变形过渡围道绘制
    \def\c{1.5}
    \def\R{-3.8}
    \def\T{2.0}
    \def\eps{0.3}
    
    \draw[thick, purple, postaction={decorate}] (\c, \eps) -- (\c, \T);
    \draw[thick, purple, postaction={decorate}] (\c, \T) -- (\R, \T);
    \draw[thick, purple, postaction={decorate}] (\R, \T) -- (\R, -\T);
    \draw[thick, purple, postaction={decorate}] (\R, -\T) -- (\c, -\T);
    \draw[thick, purple, postaction={decorate}] (\c, -\T) -- (\c, -\eps);
    
    % 槽口连线
    \draw[thick, purple, postaction={decorate}] (\c, -\eps) -- (\eps*1.732, -\eps);
    \draw[thick, purple, postaction={decorate}] (\eps*1.732, -\eps) arc (210:150:\eps);
    \draw[thick, purple, postaction={decorate}] (\eps*1.732, \eps) -- (\c, \eps);
    
    % 标注极点
    \filldraw[black] (0, 0.7) circle (1.5pt) node[left] {\small $2\pi i$};
    \filldraw[black] (0, -0.7) circle (1.5pt) node[left] {\small $-2\pi i$};
    \filldraw[black] (0, 1.4) circle (1.5pt) node[left] {\small $4\pi i$};
    \filldraw[black] (0, -1.4) circle (1.5pt) node[left] {\small $-4\pi i$};
    
    \node[purple] at (-2.0, 1.5) {\small 变形的闭合割口围道 $\mathcal{R}_{\text{cut}}$};
\end{tikzpicture}
\caption{将矩形围道沿分支割线“内折”变形为锁孔状闭合围道的几何过渡原理}
\label{fig:deformed_contour}
\end{figure}

如图 \ref{fig:deformed_contour} 所示，我们将一个顶点为 $c \pm iT, -R \pm iT$ 的矩形围道，在右侧边界 $\operatorname{Re}(s)=c$ 处剪开一个宽度为 $2\epsilon$ 的水平槽口，并将槽口沿实轴两侧“向内折叠”，环绕原点做半径为 $\epsilon$ 的半圆。这个变形后的闭合割口围道记为 $\mathcal{R}_{\text{cut}}$。
根据留数定理，若被积函数在割线 $\mathbb{C} \setminus [0,\infty)$ 外的极点集合为 $\{z_k\}$，则：
\begin{equation}
    \oint_{\mathcal{R}_{\text{cut}}} f(z) dz = 2\pi i \sum_{\text{enclosed}} \operatorname{Res}(f, z_k)
\end{equation}

现在，当我们将矩形的外侧边界推向无穷（即 $R, T \to \infty$），并令槽口的半高度 $\epsilon \to 0$ 时，这个闭合围道 $\mathcal{R}_{\text{cut}}$ 就自然分裂并过渡为两个独立的积分部分：
1. 包围整个切割线并延伸至 $+\infty$ 的\textbf{标准锁孔形围道 $\mathcal{H}$}（带有相反的绕行方向）；
2. 矩形围道的原始四边积分。

这便是两种围道在拓扑与极限层面的统一桥梁。

---

\section{积分边界贡献的严格极限估计}

要保证上述几何变形在分析上完全严格，我们必须证明：\textbf{在极限过程中，外侧边界以及槽口处的“过渡边缘”积分贡献在极限下确实消失（趋于零）}。下面我们针对本书中的两个核心应用给出详细的估算。

\subsection{锁孔形围道小圆弧极限（原点处的极性估计）}

在第\ref{ch:riemann_zeta_continuation}章 $\zeta(s)$ 的解析延拓第一法中，我们考虑了围绕原点的半径为 $\epsilon$ 的小圆弧积分 $\mathcal{C}_\epsilon$。被积函数为：
\[
  f(z) = \frac{(-z)^{s-1}}{e^z - 1}
\]
当 $\operatorname{Re}(s) = \sigma > 1$ 且 $\epsilon \to 0$ 时：
* 分子模长：$|(-z)^{s-1}| = \epsilon^{\sigma - 1} e^{-t \arg(-z)} \le \epsilon^{\sigma - 1} e^{\pi |t|}$ （其中 $s = \sigma + it$）；
* 分母渐近：$|e^z - 1| \sim |z| = \epsilon$；
* 路径长度：$L(\mathcal{C}_\epsilon) \le 2\pi \epsilon$。

因此，小圆弧上的积分贡献估算为：
\begin{equation}
    \left| \int_{\mathcal{C}_\epsilon} \frac{(-z)^{s-1}}{e^z - 1} dz \right| \le \max_{z\in\mathcal{C}_\epsilon} |f(z)| \cdot L(\mathcal{C}_\epsilon) \le C(s) \cdot \epsilon^{\sigma - 2} \cdot 2\pi \epsilon = O(\epsilon^{\sigma - 1})
\end{equation}
由于 $\sigma > 1$，当 $\epsilon \to 0$ 时，$\epsilon^{\sigma - 1} \to 0$。这严格证明了\textbf{锁孔形围道在原点处的槽口圆弧贡献确实消失}。

\subsection{矩形围道左侧垂直边极限（指数级衰减估计）}

在第\ref{ch:riemann_explicit_formula}章黎曼显式公式的 Perron 围道积分平移中，我们将垂直积分线 $\operatorname{Re}(s)=c$ 向左平移至 $\operatorname{Re}(s)=-R$。考虑在左侧边 $\mathcal{L}_L$ 上的被积函数项：
\[
  f(s) = F(s) \frac{x^s}{s}
\]
其中 $x > 1$ 是实数，$s = -R + iy$（$-T \le y \le T$）。
* 模长估计：$|x^s| = x^{-R}$；
* 分母估计：$|s| = \sqrt{R^2 + y^2} \ge R$；
* Dirichlet 级数 $F(s)$ 在左半平面的增长受限于其解析延拓后的多项式增长（由函数方程确定）。具体地，由 Phragmén--Lindelöf 凸性定理，当 $\operatorname{Re}(s) = -R \to -\infty$ 时，对有界的 $y$，有 $|F(-R+iy)| = O(R^K)$ （其中 $K$ 为常数）。

因此，左侧垂直边积分估计为：
\begin{equation}
    \left| \int_{-R-iT}^{-R+iT} F(s) \frac{x^s}{s} ds \right| \le \int_{-T}^{T} C R^K \frac{x^{-R}}{R} dy = 2T C R^{K-1} x^{-R}
\end{equation}
由于 $x > 1$，当 $R \to \infty$ 时，$x^{-R} = e^{-R \log x}$ 呈现\textbf{指数级快速衰减}。因此：
\begin{equation}
    \lim_{R\to\infty} \left| \int_{-R-iT}^{-R+iT} F(s) \frac{x^s}{s} ds \right| = 0
\end{equation}
这严格证明了\textbf{矩形围道在向左无限平移时，左边界的积分贡献完全趋于零}。

\subsection{上下水平边极限（Phragmén--Lindelöf 凸性与高度估计）}

最后，我们需要保证当高度 $T \to \infty$ 时，水平边 $\mathcal{L}_H$（$\operatorname{Im}(s) = \pm T$）的贡献也为零。以第\ref{ch:riemann_explicit_formula}章 Perron 积分估计为例，水平边位于 $s = \sigma \pm iT$（$-R \le \sigma \le c$）。
利用临界带内 $\zeta(s)$ 的估计，在垂直方向上，当 $T \to \infty$ 时，$\zeta(\sigma \pm iT)$ 仅呈多项式级缓慢增长：
\[
  |\zeta(\sigma \pm iT)| \le C T^{\frac{1-\sigma}{2}}
\]
而分母 $|s| \ge T$ 提供了 $O(T^{-1})$ 的压制。结合这两者，可以通过精细的截断估计（选取不穿过零点的 $T$），使水平边的积分贡献受限于 $O(T^{-1} \log T)$ 或更低量级，在 $T \to \infty$ 时严格趋于零（详见第\ref{ch:riemann_explicit_formula}章关于误差项的具体有界性论证）。

---

\section{结论}

综上所述，锁孔形（Hankel）围道与矩形围道并不是彼此孤立的复分析技巧，而是\textbf{在 Cauchy 积分定理下通过割口同伦变形（Homotopy）相互联系的统一体}。
* 在处理\textbf{多值分支切割}的解析延拓时，锁孔形围道能够将割线无缝包裹，是推导函数方程的最佳工具；
* 在处理\textbf{离散奇点（零点与极点）计数与部分和逼近}时，矩形围道能够清晰分离奇点，是推导显式公式的不二选择。

通过本附录的极限分析，读者可以清晰地认识到，无论选择哪种围道进行推导，它们最终捕获的奇点信息（留数值）是完全一致的。这种几何与分析上的和谐一致，正是复变函数理论高度成熟与优美的具体体现。


